% Generated draft: requires mathematical and source review.
\documentclass[11pt]{article}
\usepackage[margin=25mm]{geometry}
\usepackage{fontspec}
\setmainfont{Latin Modern Roman}
\usepackage{amsmath,amssymb,mathtools,unicode-math}
\setmathfont{Latin Modern Math}
\usepackage{enumitem,needspace,xurl,hyperref}
\setlength{\parindent}{0pt}
\setlength{\parskip}{5pt}
\setlength{\emergencystretch}{3em}
\newcommand{\N}{\mathbb{N}}
\newcommand{\Z}{\mathbb{Z}}
\newcommand{\Q}{\mathbb{Q}}
\newcommand{\R}{\mathbb{R}}
\newcommand{\C}{\mathbb{C}}
\newcommand{\F}{\mathbb{F}}
\newcommand{\A}{\mathbb{A}}
\newcommand{\PP}{\mathbb{P}}
\newcommand{\E}{\mathbb{E}}
\newcommand{\eps}{\varepsilon}
\newcommand{\dd}{\,\mathrm{d}}
\begin{document}
\section*{UnsolvedMath ID: 7000013}
\textbf{Source problem code:} catalogue entry
\subsection{Definitions and mathematical statement}
Let \(P,P'\subset\mathbb R^3\) be polyhedral surfaces, with planar polygonal faces \(\mathcal F(P)\) and \(\mathcal F(P')\). On the natural face-correspondence reading, suppose there is a bijection \(\sigma:\mathcal F(P)\to\mathcal F(P')\) such that
\[
 \operatorname{plane}(F)\parallel\operatorname{plane}(\sigma(F)),\qquad
 \operatorname{area}(F)=\operatorname{area}(\sigma(F))
 \quad(F\in\mathcal F(P)).
\]
Must \(P\) and \(P'\) be congruent, meaning that \(P'=QP+b\) for some \(Q\in O(3)\) and \(b\in\mathbb R^3\)? Reflections are allowed in congruence; no conclusion specifically of translation is imposed.

The source does not define the polyhedral-surface class: closedness, compactness, connectedness, embeddedness versus immersion, convexity, and prescribed face incidence are unspecified. It also does not explicitly give the face correspondence or require it to preserve adjacency. These are material hypotheses, not details to fill in by guessing.

Parallel face planes by themselves determine normals only up to sign. Equal oriented face-area vectors \(\operatorname{area}(F)n_F\), with consistently chosen outward normals in the convex case, give stronger data than the unoriented parallelism written above. The supplied summary invokes a convex Minkowski-type uniqueness theorem for oriented area-normal data and calls the general nonconvex case open. That theorem must not be applied to a weaker unoriented hypothesis without clarification. Its references to possible immersed counterexamples are speculative and contain no example or proof.

Thus the displayed implication retains a conventional reading of the stated matching condition, while the original geometric class and orientation convention require source review. The record's open assessment is not an independently verified current status.
\subsection{Short English statement}
Do matching parallel faces with equal areas determine a polyhedral surface up to Euclidean congruence? The surface class, face correspondence, and oriented-versus-unoriented normal data are not specified; convex oriented-area uniqueness does not automatically settle the stated general question.
\subsection{Sources}
\begin{itemize}
\item Local registry: \url{data/unsolvedmath/problems.json}, permanent UnsolvedMath ID \(7000013\), source code catalogue entry; supplied statement, background, and research summary, updated 2026-08-06.
\item Mohammad Ghomi, \emph{Open Problems in Geometry of Curves and Surfaces} (2019), Problem 2.4, PDF page 10: \url{https://people.math.gatech.edu/~ghomi/Papers/op.pdf}. Supplied normalized PDF-text extraction accessed 2026-07-29; current status marked as needing review.
\item The summary names a ``Minkowski/Robbin'' convex uniqueness theorem without a precise citation or statement. No general nonconvex counterexample or resolution reference is supplied.
\end{itemize}
\end{document}
